\documentclass[12pt]{article}
\usepackage{amsmath}
\usepackage{amssymb}
\usepackage{geometry}
\usepackage{enumerate}
\usepackage[utf8]{inputenc}
\usepackage{CJKutf8}

\geometry{margin=1in}

\title{GRE数学核心知识点总结 \\ GRE Math Core Knowledge Summary}
\author{数学笔记 \\ Math Notes}
\date{\today}

\begin{document}
\begin{CJK*}{UTF8}{gbsn}

\maketitle

\section{基础数学概念}

\subsection{因子与完全平方数}
\begin{enumerate}
    \item 正整数$n$有奇数个因子，则$n$为完全平方数
    \item 因子个数求解公式：将整数$n$分解为质因子乘积形式，然后将每个质因子的幂分别加一相乘。
    
    设$n = a^3 \times b^2 \times c$，则因子个数$= (3+1)(2+1)(1+1) = 24$个
    
    例如：$200 = 2^3 \times 5^2$，因子个数$= (3+1)(2+1) = 12$个
\end{enumerate}

\subsection{整除性判断}
\begin{enumerate}
    \item 能被8整除的数：后三位的和能被8整除
    \item 能被9整除的数：各位数的和能被9整除
    \item 能被3整除的数：各位的和能被3整除
\end{enumerate}

\subsection{几何公式}
\begin{enumerate}
    \item 多边形内角和：$(n-2) \times 180°$
    \item 菱形面积：$\frac{1}{2} \times$ 对角线乘积
    \item 欧拉公式：边数 $=$ 面数 $+$ 顶点数 $-2$
\end{enumerate}

\subsection{三角函数}
\begin{enumerate}
    \item 余弦定理：$c^2 = a^2 + b^2 - 2ab\cos\beta$，其中$\beta$为$a$、$b$两条线间的夹角
    \item 正弦定理：$\frac{a}{\sin A} = \frac{b}{\sin B} = \frac{c}{\sin C} = 2R$，其中$A$、$B$、$C$是各边及所对应的角，$R$是三角形外接圆的半径
\end{enumerate}

\subsection{直线与斜率}
\begin{enumerate}
    \item 两直线垂直条件：$y = k_1x + b_1$，$y = k_2x + b_2$，两线垂直的条件为$k_1k_2 = -1$
\end{enumerate}

\subsection{阶乘}
\begin{enumerate}
    \item $N$的阶乘公式：$N! = 1 \times 2 \times 3 \times \cdots \times (N-2) \times (N-1) \times N$
    \item 规定：$0! = 1$，$1! = 1$
    \item 例如：$8! = 1 \times 2 \times 3 \times 4 \times 5 \times 6 \times 7 \times 8$
\end{enumerate}

\subsection{常用数值}
\begin{enumerate}
    \item $\sqrt{2} = 1.414$
    \item $\sqrt{3} = 1.732$
    \item $\sqrt{5} = 2.236$
\end{enumerate}

\subsection{数量关系表达}
\begin{enumerate}
    \item $\frac{2}{3}$ as many A as B：$A = \frac{2}{3}B$
    \item twice as many A as B：$A = 2B$
\end{enumerate}

\subsection{温度换算}
\begin{enumerate}
    \item 华氏温度与摄氏温度换算：$(F-32) \times \frac{5}{9} = C$
\end{enumerate}

\subsection{常用数学公式}
\begin{enumerate}
    \item \textbf{平方差公式和完全平方公式}：
    \begin{align*}
        (a+b)(a-b) &= a^2 - b^2 \\
        (a+b)^2 &= a^2 + 2ab + b^2 \\
        (a-b)^2 &= a^2 - 2ab + b^2
    \end{align*}
    
    \item \textbf{立方和与立方差公式}：
    \begin{align*}
        (a+b)^3 &= a^3 + 3a^2b + 3ab^2 + b^3 \\
        (a-b)^3 &= a^3 - 3a^2b + 3ab^2 - b^3
    \end{align*}
    
    \item \textbf{一元二次方程求根公式}：
    
    对于方程$ax^2 + bx + c = 0$，其解为：
    $$x_{1,2} = \frac{-b \pm \sqrt{b^2 - 4ac}}{2a}$$
    
    \item \textbf{利息计算}：
    \begin{enumerate}
        \item Simple Interest（单利）：$I = P \times R \times T$
        \begin{itemize}
            \item $I$：利息（Interest）
            \item $P$：本金（Principal）
            \item $T$：时间（Time）
            \item $R$：利率（Rate）
        \end{itemize}
        \item Compound Interest（复利）：$A = P(1+R)^n$
        \begin{itemize}
            \item $A$：本利和
            \item $P$：本金
            \item $R$：利率
            \item $n$：期数
        \end{itemize}
    \end{enumerate}
    
    \item \textbf{其他常用公式}：
    \begin{align*}
        \text{Discount} &= \text{Cost} \times \text{Rate of Discount} \\
        \text{Distance} &= \text{Speed} \times \text{Time}
    \end{align*}
    
    \item \textbf{勾股定理（Pythagorean Theorem）}：
    
    直角三角形（right triangle）两直角边（legs）的平方和等于斜边（hypotenuse）的平方：
    $$a^2 + b^2 = c^2$$
    
    \item \textbf{多边形性质}：
    \begin{align*}
        \text{内角和} &= (n-2) \times 180° \\
        \text{总对角线数} &= \frac{n(n-3)}{2} \text{条} \\
        \text{每个顶点引出的对角线数} &= n-3 \text{条}
    \end{align*}
    其中$n$为多边形的边数。
    
    \item \textbf{立体图形的表面积和体积}：
    \begin{center}
    \begin{tabular}{|c|c|c|}
    \hline
    图形 & 体积（Volume） & 表面积（Surface Area） \\
    \hline
    长方体（Rectangular Prism） & 长×宽×高 & $2(\text{长} \times \text{宽} + \text{长} \times \text{高} + \text{宽} \times \text{高})$ \\
    \hline
    正方体（Cube） & 棱长的立方 & $6 \times \text{棱长} \times \text{棱长}$ \\
    \hline
    圆柱（Right Circular Cylinder） & $\pi r^2 h$ & $2\pi r h + 2\pi r^2$ \\
    \hline
    球体（Sphere） & $\frac{4\pi r^3}{3}$ & $4\pi r^2$ \\
    \hline
    圆锥（Right Circular Cone） & $\frac{\pi r^2 h}{3}$ & $\pi r^2 + \pi r l$（$l$为母线） \\
    \hline
    \end{tabular}
    \end{center}
    
    \item \textbf{平面直角坐标系中的公式}：
    
    设$A(x_1,y_1)$和$B(x_2,y_2)$是任意两点，$C(x,y)$是线段$AB$的中点，则：
    \begin{align*}
        x &= \frac{x_1 + x_2}{2} \\
        y &= \frac{y_1 + y_2}{2}
    \end{align*}
    
    线段$AB$两端点间的距离：
    $$d = \sqrt{(x_1-x_2)^2 + (y_1-y_2)^2}$$
\end{enumerate}

\subsection{概率论基础}
\begin{enumerate}
    \item \textbf{排列(Permutation)}：从$N$个东东(有区别)中不重复(即取完后不再取)取出$M$个并作排列，共有几种方法：
    $$P(M,N) = \frac{N!}{(N-M)!}$$
    
    例如：从1-5中取出3个数不重复，问能组成几个三位数？
    解答：$P(3,5) = \frac{5!}{(5-3)!} = \frac{5!}{2!} = \frac{5 \times 4 \times 3 \times 2 \times 1}{2 \times 1} = 5 \times 4 \times 3 = 60$
    
    也可以这样想：从五个数中取出三个放三个固定位置，第一个位置可以放五个数中任一一个，所以有5种可能选法，第二个位置余下四个数中任一个，有4种选法，第三个位置有3种选法，所以总共的排列为$5 \times 4 \times 3 = 60$。
    
    同理可知如果可以重复选(即取完后可再取)，总共的排列是$5 \times 5 \times 5 = 125$。
    
    \item \textbf{组合(Combination)}：从$N$个东东(可以无区别)中不重复(即取完后不再取)取出$M$个(不作排列，即不管取得次序先后)，共有几种方法：
    $$C(M,N) = \frac{P(M,N)}{P(M,M)} = \frac{N!}{(N-M)!M!}$$
    
    例如：$C(3,5) = \frac{P(3,5)}{P(3,3)} = \frac{5!}{2!3!} = \frac{5 \times 4 \times 3}{1 \times 2 \times 3} = 10$
    
    可以这样理解：组合与排列的区别就在于取出的$M$个作不作排列，即$M$的全排列$P(M,M) = M!$，那么他们之间关系就有先做组合再作$M$的全排列就得到了排列，所以$C(M,N) \times P(M,M) = P(M,N)$，由此可得组合公式。
    
    性质：$C(M,N) = C(N-M, N)$，即$C(3,5) = C(5-3, 5) = C(2,5) = \frac{5!}{3!2!} = 10$。
    
    \item \textbf{概率}：概率的定义：$P = \frac{\text{满足某个条件的所有可能情况数量}}{\text{所有可能情况数量}}$
    
    概率的性质：$0 \leq P \leq 1$
    
    \begin{enumerate}
        \item \textbf{不相容事件的概率}：$a$、$b$为两两不兼容的事件（即发生了$a$，就不会发生$b$）
        $$P(a \text{或} b) = P(a) + P(b)$$
        $$P(a \text{且} b) = P(a) + P(b) = 0 \quad (A,B \text{不能同时发生})$$
        
        \item \textbf{对立事件的概率}：对立事件就是$a + b$就是全部情况，所以不是发生$a$，就是$b$发生，但是，有一点$a$、$b$不能同时发生。
        
        例如：$a$：一件事不发生，$b$：一件事发生，则$A$、$B$是对立事件。
        
        显然：$P(\text{一件事发生的概率或一件事不发生的概率}) = 1$（必然事件的概率为1），则一件事发生的概率$= 1 - $一件事不发生的概率。
        
        \item \textbf{条件概率}：考虑的是事件$A$已发生的条件下事件$B$发生的概率。
        
        定义：设$A$、$B$是两个事件，且$P(A) > 0$，称
        $$P(B|A) = \frac{P(A \cap B)}{P(A)}$$
        为事件$A$已发生的条件下事件$B$发生的概率。
        
                 \item \textbf{独立事件与概率}：两个事件独立也就是说，$A$、$B$的发生与否互不影响，$A$是$A$，$B$是$B$，用公式表示就是$P(A|B) = P(A)$，所以说两个事件同时发生的概率就是：
        $$P(A \cap B) = P(A) \times P(B)$$
    \end{enumerate}
\end{enumerate}

\subsection{统计学基础}
\begin{enumerate}
    \item \textbf{Mode（众数）}：一堆数中出现频率最高的一个或几个数。
    
    例如：mode of 1,1,1,2,3,0,0,0,5 is 1 and 0。
    
    \item \textbf{Range（值域）}：一堆数中最大和最小数之差，所以统计学上又称之为极差（两极的差）。
    
    例如：range of 1,1,2,3,5 is 5-1=4。
    
    \item \textbf{Mean（平均数）}：
    \begin{enumerate}
        \item arithmetic mean（算术平均数）：$n$个数之和再除以$n$
        \item geometric mean（几何平均数）：$n$个数之积的$n$次方根
    \end{enumerate}
    
    \item \textbf{Median（中数）}：将一堆数排序之后，正中间的一个数（奇数个数位），或者中间两个数的平均数（偶数个数字）。
    
    例如：median of 1,7,4,9,2,2,2,2,2,5,8 is 2；median of 1,7,4,9,2,5 is $(5+7)/2=6$。
    
    \item \textbf{Standard Error（标准偏差）}：一堆数中，每个数与平均数的差的绝对值之和，除以这堆数的个数($n$)。
    
    例如：standard error of 0,2,5,7,6 is：$(|0-4|+|2-4|+|5-4|+|7-4|+|6-4|)/5=2.4$。
    
    \item \textbf{Standard Variation（标准方差）}：一堆数中，每个数与平均数之差的平方之和，再除以$n$。
    
    标准方差的公式：$d^2 = \frac{(a_1-a)^2+(a_2-a)^2+\cdots+(a_n-a)^2}{n}$
    
    例如：standard variation of 0,2,5,7,6 is：average=4，$((0-4)^2+(2-4)^2+(5-4)^2+(7-4)^2+(6-4)^2)/5=6.8$。
    
    \item \textbf{Standard Deviation（标准差）}：就是standard variation的平方根$d$。
    
    \item \textbf{四分位数的计算}：
    \begin{enumerate}
        \item 第0个Quartile实际为通常所说的最小值（MINimum）
        \item 第1个Quartile（1st Quartile）
        \item 第2个Quartile实际为通常所说的中分位数（中数、二分位分、中位数：Median）
        \item 第3个Quartile（3rd Quartile）
        \item 第4个Quartile实际为通常所说的最大值（MAXimum）
    \end{enumerate}
    
    求1st Quartile的步骤：
    \begin{enumerate}
        \item $n$个数从小到大排列，求$(n-1)/4$，设商为$i$，余数为$j$
        \item 则可求得1st Quartile为：$(第i+1个数) \times (4-j)/4 + (第i+2个数) \times j/4$
    \end{enumerate}
    
    例如：
    \begin{enumerate}
        \item 序列为$\{5\}$，只有一个样本则：$(1-1)/4$ 商0，余数0，1st=第1个数$\times 4/4$+第2个数$\times 0/4=5$
        \item 序列为$\{1,4\}$，有两个样本则：$(2-1)/4$ 商0，余数1，1st=第1个数$\times 3/4$+第2个数$\times 1/4=1.75$
        \item 序列为$\{1,5,7\}$，有三个样本则：$(3-1)/4$ 商0，余数2，1st=第1个数$\times 2/4$+第2个数$\times 2/4=3$
        \item 序列为$\{1,3,6,10\}$，四个样本：$(4-1)/4$ 商0，余数2，1st=第1个数$\times 1/4$+第2个数$\times 3/4=2.5$
    \end{enumerate}
    
    因为3rd与1st的位置对称，可以将序列从大到小排（即倒过来排），再用1st的公式即可求得。
    
    \item \textbf{百分位数的计算}：
    
    设一个序列共有$n$个数，要求$(k\%)$的Percentile：
    \begin{enumerate}
        \item 从小到大排序，求$(n-1) \times k\%$，记整数部分为$i$，小数部分为$j$
        \item 所求结果$= (1-j) \times 第(i+1)个数 + j \times 第(i+2)个数$
    \end{enumerate}
    
    特别注意以下两种最可能考的情况：
    \begin{enumerate}
        \item $j$为0，即$(n-1) \times k\%$恰为整数，则结果恰为第$(i+1)$个数
        \item 第$(i+1)$个数与第$(i+2)$个数相等，不用算也知道正是这两个数
    \end{enumerate}
    
    前面提到的Quartile也可用这种方法计算，其中1st Quartile的$k\%=25\%$，2nd Quartile的$k\%=50\%$，3rd Quartile的$k\%=75\%$。
    
    例如：$\{1,3,4,5,6,7,8,9,19,29,39,49,59,69,79,80\}$共16个样本，要求percentile$=30\%$：
    $(16-1) \times 30\% = 4.5 = 4 + 0.5$，$i=4$，$j=0.5$，$(1-0.5) \times 第5个数 + 0.5 \times 第6个数 = 0.5 \times 6 + 0.5 \times 7 = 6.5$。
    
    \item \textbf{茎叶法计算中位数}：
    
    Stem-and-Leaf method其实并不是很适用于GRE考试，除非有大量数据时可以用这种方法比较迅速的将数据有序化。一般GRE给出的数据在10个左右，茎叶法有点大材小用。
    
    Stem-and-Leaf其实就是一种分级将数据分类的方法。Stem就是大的划分，如可以划分为1～10，11～20，21～30…，而Leaf就是把划分到Stem一类中的数据再排一下序。
    
    例如：Data：23，51，1，24，18，2，2，27，59，4，12，23，15，20
    \begin{verbatim}
    0|  1  2  2  4 
    1| 12 15 18 
    2| 20 23 23 24 27 
    5| 51 59 
    Stem (unit) = 10 
    Leaf (unit) = 1
    \end{verbatim}
    
    最左边的一竖行0, 1, 2, 5叫做Stem，而右边剩下的就是Leaf(leaves)。上面的Stem-and-Leaf共包含了14个data，根据Stem及leaf的unit，分别是：1, 2, 2, 4 (first row), 12, 15, 18 (second row), 20, 23, 23, 24, 27(third row), 51, 59 (last row)。
    
    \item \textbf{按比例求中位数}：
    
    给了不同年龄range，和各个range的percentage，问median落在哪个range里。把percentage加到50\%就是median的range了。但小心一点，range首先要保证是有序排列。
    
    例如：Given: 10～20 = 20\%, 30～50 = 30\%, 0～10 = 40\%, 20～30 = 10\%，问median在哪个range里。
    
    分析：千万不要上来就加，要先排序，切记！重新排序为：0～10 = 40\%, 10～20 = 20\%, 20～30 = 10\%, 30～50 = 40\%。然后从小开始加，median（50\%）落在10～20这个range里。
    
    \item \textbf{算术平均值和加权平均值}：
    
    三组资料的频数分布FREQUENCY DISTRIBUTION：
    \begin{enumerate}
        \item 1（6），2（4），3（1），4（4），5（6）
        \item 1（1），2（4），3（6），4（4），5（1）
        \item 1（1），2（2），3（3），4（4），5（5）
    \end{enumerate}
    
    其中括号里的是出现的频率，问MEAN和AVERAGE相等的有那些。
    
    答案：只有第二个。
    mean-arithmetic mean算术平均值：$(1+2+3+4+5)/5 = 3$
    average-weighted average加权平均值：$(1 \times 1+2 \times 4+\cdots+5 \times 1)/(1+4+6+4+1)=48/16=3$
    
    \item \textbf{正态分布}：
    
    高斯分布（Gaussian）（正态分布）的概率密度函数为一钟型曲线，即：
    $$f(x) = \frac{1}{\sigma \sqrt{2\pi}} e^{-\frac{(x-a)^2}{2\sigma^2}}$$
    
    其中$a$为均值，$\sigma$为标准方差，曲线关于$x=a$的虚线对称，$\sigma$决定了曲线的"胖瘦"。
    
    正态分布的分布函数$F(x)$在其期望$a$的右方曲线是向上凸的，此时如果$F(20)=75\%$，$F(r)=85\%$，$F(26)=95\%$，由于曲线上凸，显然$r < 23$。
    
    对于一般的正态分布，可以通过变换，归一化到标准的正态分布，算法为：
    设原正态分布的期望为$a$，标准方差为$\sigma$，欲求分布在区间$(y_1, y_2)$的概率，可以变换为求标准正态分布中分布在$(x_1, x_2)$间的概率，其中：
    $$x_1 = \frac{y_1 - a}{\sigma}, \quad x_2 = \frac{y_2 - a}{\sigma}$$
\end{enumerate}

\section{练习题与解答}

\subsection{数列题}
\begin{enumerate}
    \item 已知$a_1 = 2$，$a_2 = 6$，$a_n = \frac{a_{n-1}}{a_{n-2}}$，求$a_{150}$。
    
    \textbf{解答：} $a_n = \frac{a_{n-1}}{a_{n-2}}$，所以$a_{n-1} = \frac{a_{n-2}}{a_{n-3}}$，带入前式得$a_n = \frac{1}{a_{n-3}}$，然后再拆一遍得到$a_n = a_{n-6}$，也就是说，这个数列是以6为周期的，则$a_{150} = a_{144} = \cdots = a_6$，利用$a_1$、$a_2$可以计算出$a_6 = \frac{1}{3}$。
    
    如果实在想不到这个方法，可以写几项看看很快就会发现$a_{150} = a_{144}$，大胆推测该数列是以6为周期的，然后写出$a_1$到$a_{13}$（也就是写到你能看出来规律），不难发现$a_6 = a_{12}$，$a_7 = a_{13}$，然后稍微数数，就可以知道$a_{150} = a_6$了，同样计算得$\frac{1}{3}$。
\end{enumerate}

\subsection{温度换算题}
\begin{enumerate}
    \item 问摄氏升高30度华氏升高的度数与62比大小。
    
    \textbf{解答：} $F = 30 \times \frac{9}{5} = 54 < 62$
\end{enumerate}

\subsection{斐波那契数列题}
\begin{enumerate}
    \item 已知$a_1 = 1$，$a_2 = 1$，$a_n = a_{n-1} + a_{n-2}$，问$a_1$、$a_2$、$a_3$、$a_6$四项的平均数和$a_1$、$a_3$、$a_4$、$a_5$四项的平均数大小比较。
    
    \textbf{解答：} 斐波那契数列就是第三项是前两项的和，依此类推得到$a_1$到$a_6$为：1, 1, 2, 3, 5, 8, 13, 21。$a_1 + a_2 + a_3 + a_6 = 12$，$a_1 + a_3 + a_4 + a_5 = 11$，所以为大于。
\end{enumerate}

\subsection{整数对计数题}
\begin{enumerate}
    \item 满足$x^2 + y^2 \leq 100$的整数对$(x,y)$有多少？
    
    \textbf{解答：} 按照$X$的可能情况顺序写出：
    \begin{align*}
        X &= 1-9 \quad Y = 1-9 \\
        X &= 1-9 \quad Y = 1-9 \\
        X &= 1-9 \quad Y = 1-9 \\
        X &= 1-9 \quad Y = 1-9 \\
        X &= 1-8 \quad Y = 1-8 \\
        X &= 1-8 \quad Y = 1-8 \\
        X &= 1-7 \quad Y = 1-7 \\
        X &= 1-6 \quad Y = 1-6 \\
        X &= 1-4 \quad Y = 1-4
    \end{align*}
    加起来$= 69$
\end{enumerate}

\subsection{质因子题}
\begin{enumerate}
    \item 24，36，90，100四个数中，该数除以它的所有的质因子，最后的结果是质数的是那个？
    
    \textbf{解答：} 90
\end{enumerate}

\subsection{小数位数题}
\begin{enumerate}
    \item 0.123456789101112…，这个小数无限不循环地把所有整数都列出来。请问小数点后第100位的数字是多少？
    
    \textbf{解答：}
    \begin{align*}
        &1 \text{到} 9: \text{9位} \\
        &10 \text{到} 19: \text{20位} \\
        &20 \text{到} 29: \text{20位} \\
        &30 \text{到} 39: \text{20位} \\
        &40 \text{到} 49: \text{20位} \\
        &50 \text{到} 56: \text{第101位} = 5
    \end{align*}
\end{enumerate}

\subsection{完全平方数题}
\begin{enumerate}
    \item $2904x = y^2$（$y$的平方），$x$、$y$都是正整数，求$x$的最小值。
    
    \textbf{解答：} 因为$X^2 \times Y^2 \times Z^2 = (X \times Y \times Z)^2$，所以把2904除呀除$= 2 \times 2 \times 2 \times 3 \times 11 \times 11 = 2^2 \times 11^2 \times 6$，再乘一个6就OK了。$2^2 \times 11^2 \times 6 \times 6 = (2 \times 11 \times 6)^2 = 132^2$。
    
    \textbf{答案：} 最小的$x = 6$
\end{enumerate}

\subsection{数列求和题}
\begin{enumerate}
    \item 序列$a_n = \frac{1}{n} - \frac{1}{n+1}$，$n \geq 1$，问前100项和。
    
    \textbf{解答：}
    \begin{align*}
        a_n &= \frac{1}{n} - \frac{1}{n+1} \\
        a_{n-1} &= \frac{1}{n-1} - \frac{1}{n} \\
        a_{n-2} &= \frac{1}{n-2} - \frac{1}{n-1} \\
        &\vdots \\
        a_1 &= 1 - \frac{1}{2}
    \end{align*}
    把左边加起来就是$a_n + a_{n-1} + \cdots + a_1 = 1 - \frac{1}{n+1}$，消掉了好多好多项之后的结果。
    
    \textbf{答案：} 把$n = 100$带入得前100项之和为$\frac{100}{101}$
\end{enumerate}

\subsection{等腰三角形题}
\begin{enumerate}
    \item 等腰三角形，腰为6，底边上的高为$x$，底边为$y$，问$4x^2 + y^2$和144谁大。
    
    \textbf{解答：} 勾股定理得$(\frac{y}{2})^2 + x^2 = 6^2$，所以$4x^2 + y^2 = 144$
\end{enumerate}

\subsection{不等式题}
\begin{enumerate}
    \item $-1 < r < t < 0$（有一数轴），问：$r + rt^2$与$-1$的关系。
    
    \textbf{解答：} 我想的办法只能是尝试：
    原式$= r(1 + t^2)$恒小于零
    \begin{enumerate}
        \item $r \rightarrow -1$，$t \rightarrow 0$，则原式$\rightarrow -1$
        \item $r \rightarrow -1$，$t \rightarrow -1$，则原式$\rightarrow -2$
        \item $r \rightarrow 0$，$t \rightarrow 0$，则原式$\rightarrow 0$
    \end{enumerate}
    例如：$r = -0.9$，$t = -\frac{1}{3}$时，原式$= -1$，若此时$-0.9 < t < -\frac{1}{3}$，原式$< -1$，反之$> -1$。
\end{enumerate}

\subsection{长方形截取题}
\begin{enumerate}
    \item 有长方形4feet $\times$ 8feet，长宽各截去$x$ inch，长宽比2：5。
    
    \textbf{解答：} 列出方程：$\frac{4 \times 12 - x}{8 \times 12 - x} = \frac{2}{5}$，解得$x = 16$
\end{enumerate}

\subsection{概率论练习题}
\begin{enumerate}
    \item $A$、$B$独立事件，一个发生的概率是0.6，一个是0.8，问：两个中发生一个或都发生的概率？
    
    \textbf{解答：}
    $$P = P(A \text{且} \overline{B}) + P(B \text{且} \overline{A}) + P(A \text{且} B)$$
    $$= 0.6 \times (1-0.8) + 0.8 \times (1-0.6) + 0.6 \times 0.8 = 0.92$$
    
    另一个角度，所求概率$P = 1 - P(A,B \text{都不发生})$
    $$= 1 - (1-0.8) \times (1-0.6) = 0.92$$
    
    \item 一道概率题：就是100以内取两个数是6的整倍数的概率。
    
    \textbf{解答：} 100以内的倍数有6, 12, 18, ..., 96共计16个，所以从中取出两个共有$16 \times 15$种方法，从1-100中取出两个数的方法有$99 \times 100$种，所以$P = \frac{16 \times 15}{99 \times 100} = \frac{12}{505} = 0.024$。
    
    \item 1-350 inclusive中，在100-299 inclusive之间以3,4,5,6,7,8,9结尾的数的概率。
    
    因为100-299中以3,4,5,6,7,8,9结尾的数各有20个，所以
    \textbf{答案：} $\frac{2 \times 10 \times 7}{350} = 0.4$
    
    \item 在1-350中（inclusive），337-350之间整数占的百分比。
    
    \textbf{答案：} $\frac{350-337+1}{350} = 4\%$
    
    \item 在$E$发生的情况下，$F$发生的概率为0.45，问$E$不发生的情况下，$F$发生的概率与0.55比大小。
    
    \textbf{解答：} 用全概率公式解决这个问题。某一个事件$A$的发生总是在一定的其他条件下如$B$、$C$、$D$发生的，也就是说$A$的概率其实就是在$B$、$C$、$D$发生的条件下$A$发生的概率之和。
    
    $P(A) = P(A|B) + P(A|C) + P(A|D)$
    
    好了，看看这个题目就明白了。$F$发生时，$E$要么发生，要么不发生，所以：
    $$P(F) = P(F|E) + P(F|\overline{E})$$
    
    给了$P(F|E) = 0.45$，所以：
    $$P(F|\overline{E}) = P(F) - P(F|E) = P(F) - 0.45$$
    
    如果$P(F) = 1$，那么$P(F|\overline{E}) = 0.55$
    如果$0.45 \leq P(F) < 1$，那么$0 \leq P(F|\overline{E}) < 0.55$
\end{enumerate}

\subsection{统计学练习题}
\begin{enumerate}
    \item 比较，当$n < 1$时，$n$，1，2和1，2，3的标准方差谁大。
    
    \textbf{解答：} standard error和standard variation（作用=standard deviation）都是用来衡量一组资料的离散程度的统计数值，只不过由于standard error中涉及绝对值，在数学上是很难处理的，所以都用标准方差。实际上standard error更合理一些，它代表了数据和平均值的平均距离。很明显题目中如果$n=0$的话，0,1,2的离散程度应该和1,2,3的离散程度相同。如果$n<0$，则$n$,1,2的离散程度大于后者，而$0<n<1$的话，则后者大于前者，但是$n$为整数，这种情况不成立。故而：
    
    \textbf{答案：} $n$是整数，前$\geq$后（$n=0$，等；$n=-1,-2,\ldots$大于）。
    
    \item 正态分布题：一列数从0到28，给出正态分布曲线。75\%的percentile是20，85\%的percentile是$r$，95\%的percentile是26，问$r$与23的大小。
    
    \textbf{答案：} $r < 23$
    
    \textbf{解答：} 由图2，正态分布的分布函数$F(x)$在其期望$a$的右方曲线是向上凸的，此时$F(20)=75\%$，$F(r)=85\%$，$F(26)=95\%$。如果把曲线的片段放大就比较清楚了。$O$为$AB$的中点。
    
    $A(20, 75\%)$，$B(26, 95\%)$，$O(23, 85\%)$，$C(r, 85\%)$
    
    由于曲线上凸，显然$C$的横坐标小于$O$，所以$r < 23$。
    
    \item 正态分布题好象是：有一组数平均值9，标准方差2，另一组数平均值3，标准方差1，问分别在（5，11）和（1，4）中个数（概率）谁大，应该是相等。
    
    \textbf{解答：} 令图1中的曲线$a=0$，$\sigma=1$，就得到了标准正态分布。对于一般的正态分布，可以通过变换，归一化到标准的正态分布，算法为：
    
    设原正态分布的期望为$a$，标准方差为$\sigma$，欲求分布在区间$(y_1, y_2)$的概率，可以变换为求标准正态分布中分布在$(x_1, x_2)$间的概率，其中：
    $$x_1 = \frac{y_1 - a}{\sigma}, \quad x_2 = \frac{y_2 - a}{\sigma}$$
    
    比如题目中$a=9$，$\sigma=2$，区间为（5, 11），则区间归一化为$(-2, 1)$。
    
    同理，$a=3$，$\sigma=1$，区间为（1, 4），则区间归一化后也为$(-2, 1)$。
    
    所以两者的分布概率相等。
    
    \item 原题：说一堆人0-10岁占10\%，11-20岁占12\%，21-30岁占23\%，31-40岁占20\%，>40岁占35\%，问median在什么范围？
    
    \textbf{解答：} 按比例求中位数，把percentage加到50\%就是median的range了。但小心一点，range首先要保证是有序排列。
    
    重新排序为：0～10 = 10\%，11～20 = 12\%，21～30 = 23\%，31～40 = 20\%，>40 = 35\%
    
    然后从小开始加：10\% + 12\% + 23\% = 45\% < 50\%，10\% + 12\% + 23\% + 20\% = 65\% > 50\%
    
    所以median落在31～40这个range里。
    
    \textbf{答案：} 31～40岁
\end{enumerate}

\end{CJK*}
\end{document}
